% Part: sets-functions-relations
% Chapter: arithmetization
% Section: cuts
%
\documentclass[../../../include/open-logic-section]{subfiles}

\begin{document}

\olfileid[fr]{sfr}{arith}{cuts}
\olsection{De $\Rat$ à $\Real$}

La propriété de la borne supérieure exprime, pour l'essentiel, la manière
dont un point $\alpha$ partage la droite réelle sans laisser de lacune :
il est la borne supérieure des points situés à sa gauche et la borne
inférieure de ceux situés à sa droite. Pour \emph{construire} les nombres
réels à partir des rationnels, Dedekind a proposé de prendre pour réels
les \emph{coupures} qui partagent les rationnels. Nous identifions ainsi
$\sqrt{2}$ à la \emph{coupure} qui sépare les rationnels
$< \sqrt{2}$ des rationnels $> \sqrt{2}$.

Précisons cette idée. Lorsque nous partageons les rationnels en deux
parties, la partie \emph{inférieure} suffit à déterminer la partition.
Nous adoptons donc la définition suivante :

\begin{defn}[Coupure]
Une \emph{coupure} $\alpha$ est un segment initial non vide et propre
de l'ensemble des rationnels, sans plus grand élément. Autrement dit,
$\alpha$ est une coupure si et seulement si :
\begin{enumerate}
	\item \emph{non vide et propre} : $\emptyset \neq \alpha \subsetneq \Rat$ ;
	\item \emph{segment initial} : pour tous $p,q \in \Rat$, si $p < q \in \alpha$, alors $p \in \alpha$ ;
	\item \emph{sans plus grand élément} : pour tout $p \in \alpha$, il existe $q \in \alpha$ tel que $p < q$.
\end{enumerate}
Nous définissons alors $\Real$ comme l'ensemble des coupures.
\end{defn}

Nous pouvons maintenant poser $\sqrt{2} = \Setabs{p \in \Rat}{p^2 < 2\text{
ou }p < 0}$. Bien entendu, il faut vérifier qu'il s'agit \emph{bien}
d'une coupure ; nous reportons cette vérification à la \olref[check]{sec}.

Comme précédemment, après avoir défini les objets, nous devons définir
les fonctions et relations usuelles sur leur ensemble. Commençons par
une relation simple :
\begin{align*}
	\alpha \leq \beta \liff \alpha \subseteq \beta.
\end{align*}
Cette définition de l'ordre permet d'\emph{énoncer} le résultat central :
l'ensemble des coupures possède la propriété de la borne supérieure.
Explicitons l'énoncé. Si $S$ est un ensemble non vide de coupures qui
admet un majorant, alors il admet un plus petit majorant. Plus précisément,
il existe une coupure $\lambda$ qui majore $S$, c'est-à-dire que
$(\forall \alpha \in S)\alpha \subseteq \lambda$, et $\lambda$ est la plus
petite de ces coupures, c'est-à-dire que
$(\forall \beta \in \Real)((\forall \alpha \in S)\alpha \subseteq \beta \lif \lambda \subseteq \beta)$.
Voici la démonstration :

\begin{thm}\ollabel{realcompleteness}
L'ensemble des coupures possède la propriété de la borne supérieure.
\end{thm}

\begin{proof}
Soit $S$ un ensemble non vide de coupures admettant un majorant. Posons
$\lambda = \bigcup S$.
%\Setabs{p \in \Rat}{(\exists \alpha \in S)p \in \alpha}$$
Montrons d'abord que $\lambda$ est une coupure :
\begin{enumerate}
\item Puisque $S$ est non vide, il contient au moins une coupure, elle-même
non vide ; ainsi $\emptyset \neq \lambda$.\footnote{Correction éditoriale : la source déduit à tort l'existence d'une coupure dans $S$ du fait que $S$ admet un majorant. C'est l'hypothèse $S\neq\emptyset$ qui intervient ici. La majoration est utilisée ensuite pour montrer que la réunion est propre.}
Puisque $S$ est un ensemble de coupures, $\lambda \subseteq \Rat$.
Comme $S$ admet une coupure majorante, un rationnel $p \in \Rat$
extérieur à cette coupure n'appartient à aucune coupure $\alpha \in S$.
Ainsi $p\notin \lambda$, donc $\lambda \subsetneq \Rat$.
\item Supposons que $p < q \in \lambda$. Il existe alors $\alpha \in S$
tel que $q \in \alpha$. Comme $\alpha$ est une coupure, $p \in \alpha$,
donc $p \in \lambda$.
\item Supposons que $p \in \lambda$. Il existe alors $\alpha \in S$
tel que $p \in \alpha$. Comme $\alpha$ est une coupure, il existe
$q \in \alpha$ tel que $p < q$. Nous avons donc aussi $q \in \lambda$.
\end{enumerate}
L'affirmation est démontrée. De plus,
$(\forall \alpha \in S)\alpha \subseteq \bigcup S = \lambda$ :
$\lambda$ est donc un majorant de $S$. Supposons maintenant que
$\beta \in \mathbb{R}$ soit un autre majorant, c'est-à-dire que
$(\forall \alpha \in S)\alpha \subseteq \beta$. Pour tout $p \in \Rat$,
si $p \in \lambda$, il existe $\alpha \in S$ tel que $p \in \alpha$,
et donc $p \in \beta$. Nous en déduisons $\lambda \subseteq \beta$.
Ainsi, $\lambda$ est le \emph{plus petit} majorant de $S$.
\end{proof}

Nous avons donc construit des objets qui satisfont la propriété de la
borne supérieure. Nous pouvons exprimer ce résultat en disant que nos
coupures ne présentent aucune «~lacune~» : former de nouvelles
«~coupures~» de réels, plutôt que de rationnels, ne fournirait pas
de nouveaux objets mathématiquement intéressants.

Définissons maintenant les opérations sur les réels. Commençons par
plonger les rationnels dans les réels en posant
$p_\Real = \Setabs{q \in \Rat}{q < p}$ pour chaque $p \in \Rat$.
Définissons ensuite :
\begin{align*}
	\alpha + \beta &= \Setabs{p + q}{p \in \alpha \land q \in \beta}\\
%	\alpha - \beta &\defis \Setabs{p - q}{p \in \alpha \land q \in \Rat \setminus \beta}\\
	\alpha \times \beta &=
	\Setabs{p \times q}{0 \leq p \in \alpha \land 0 \leq q \in \beta} \cup 0_\Real & \text{si }\alpha, \beta \geq 0_\Real.
%	\alpha \div \beta &\defis
%	\Setabs{\nicefrac{p}{q}}{p \in \alpha \land q \in \Rat \setminus \beta}  & \text{si }\alpha \geq 0_\Real, \beta > 0_\Real
\end{align*}
Pour traiter les autres cas du produit, posons d'abord
$0_\Real \times \alpha = 0_\Real = \alpha \times 0_\Real$ pour toute
coupure $\alpha$, puis définissons l'opposé :
% La source suggère en commentaire de poser d'abord l'annulation par la coupure nulle ; cette règle est rendue explicite ci-dessus.
\begin{align*}
	-\alpha &= \Setabs{p - q}{p,q \in \Rat \land p < 0 \land q \notin \alpha}.
\end{align*}
Posons enfin :
\begin{align*}
	\alpha \times \beta &\defis
	\begin{cases}
		\mathord{-}\alpha \times \mathord{-}\beta &\text{si }\alpha < 0_\Real\text{ et }\beta < 0_\Real\\
		\mathord{-}(\mathord{-}\alpha \times \beta) &\text{si }\alpha < 0_\Real \text{ et }\beta > 0_\Real\\
		\mathord{-}(\alpha \times \mathord{-}\beta) &\text{si }\alpha > 0_\Real \text{ et }\beta < 0_\Real.
	\end{cases}
\end{align*}
Il faut vérifier que chacune de ces définitions produit toujours une
coupure. Il restera enfin à établir, par une démonstration élémentaire
mais longue, que les coupures ainsi munies de ces opérations ont toutes
les propriétés attendues. Nous reportons ces vérifications à la
\olref[check]{sec}.\footnote{Corrections et précisions éditoriales : la notation $0^\mathbb{R}$ du produit positif est remplacée par $0_\Real$, la coupure nulle déjà définie. Les cas où un facteur est nul et l'autre strictement négatif manquaient dans les définitions actives ; la règle d'annulation, suggérée par un commentaire source, les complète. Dans la définition de l'opposé, les domaines rationnels de $p$ et $q$, implicites dans le contexte original, sont écrits explicitement.}
\end{document}
